% !TeX root = ../初等数学笔记.tex
% 由主文档通过 \input 载入；请编译主文档。
\setcounter{chapter}{3} % 保留原书章号。
\chapter{代数式与多项式}
\label{chap:em-expressions}
代数式把具体数量关系写成可重复使用的结构。一个可靠的变形既要保持数值关系，也要保留变量范围；因式分解、配方和换元的目的，是使表达式更适合当前的求解或证明任务。

\section{变量、代数式与恒等关系}
变量表示某个给定范围内可变的数，参数在一次求解中视为固定，但可用于比较一族问题。\emph{恒等式}（identity）是在所声明范围内对每个变量取值都成立的等式，例如 $(x+1)^2=x^2+2x+1$ 对一切实数成立；$x^2=1$ 则只对特定值成立。

表达式的定义域是使其全部运算合法的取值集合。分式要求分母非零，偶次根式要求被开方数非负。$\frac{x^2-1}{x-1}=x+1$ 在 $x\ne1$ 时成立；右端虽然在 $x=1$ 也有值，却不能据此扩大原式的定义域。这一区分在连续延拓与解方程时尤其重要。

\section{多项式的运算与恒等式}
实系数一元\emph{多项式}（polynomial）写为 $f(x)=a_0+a_1x+\cdots+a_nx^n$。若 $a_n\ne0$，则次数为 $n$、首项系数为 $a_n$；零多项式不赋予通常的非负整数次数。由分配律可得
\[
 (a\pm b)^2=a^2\pm2ab+b^2,\quad
 a^2-b^2=(a-b)(a+b),\quad
 a^3\pm b^3=(a\pm b)(a^2\mp ab+b^2).
\]
这些公式都是运算律的结果，而非额外假设。非零多项式的乘积次数等于次数之和，因为最高次项系数相乘非零；和的次数却可能因首项抵消而降低。

多项式作为形式表达式相等，要求每一幂次的系数相同。把它们视为实数函数时，在所有实数上取值相同也能推出系数相同，依据是本章后面证明的根数界；仅比较几个数值一般不足以认定恒等。

\section{因式分解的结构方法}
分解先寻找公共因子，再观察平方差、完全平方或可分组的结构。例如
\[
 x^3+x^2-x-1=(x+1)(x^2-1)=(x-1)(x+1)^2.
\]
配方把二次式转化为平方与常量的组合；换元则把反复出现的子表达式作为整体，例如令 $t=x^2$ 可把 $x^4-5x^2+4$ 分解为 $(t-1)(t-4)$，再回代得到 $(x-1)(x+1)(x-2)(x+2)$。

分解依赖系数范围。$x^2-2$ 在有理数范围内没有一次因子，否则有理根将等于 $\pm\sqrt2$；在实数范围却可写成 $(x-\sqrt2)(x+\sqrt2)$。$x^2+1$ 在实数中无根，不能分解成实系数一次因子的乘积；复数会改变这一结论。

\section{多项式除法、余式与因式}
\begin{theorem}[多项式带余除法]
对实系数多项式 $f,g$ 且 $g\ne0$，存在唯一多项式 $q,r$ 使 $f=qg+r$，其中 $r=0$ 或 $\deg r<\deg g$。
\end{theorem}
\begin{proof}
当被处理多项式次数不小于 $\deg g$ 时，减去适当单项式倍的 $g$ 消去最高次项。每次次数严格下降，有限步后得到所需余式。若有两个表示，则 $(q-q')g=r'-r$；若 $q-q'\ne0$，左端次数至少为 $\deg g$，右端却小于它，矛盾。
\end{proof}
取 $g(x)=x-a$，余式为常数，代入 $x=a$ 得余式为 $f(a)$。于是\emph{因式定理}（factor theorem）为
\begin{equation}\label{eq:em-factor}
 f(a)=0\Longleftrightarrow f(x)=(x-a)q(x).
\end{equation}
非零 $n$ 次多项式至多有 $n$ 个不同实根：$n=0$ 时没有根；若 $a$ 为一根，除去因子 $x-a$ 后，其余不同根都是次数 $n-1$ 的商的根，归纳完成。因此两个次数至多为 $n$ 的多项式，若在 $n+1$ 个不同点取值相同，其差必为零多项式。

\section{分式、根式与等价化简}
分式加减先通分，乘除先检查因子与除数是否为零。例如
\[
 \frac1{x-1}-\frac1{x+1}=\frac2{x^2-1},\qquad x\ne\pm1.
\]
约分只约乘积中的公共因子，不能把和中的一项随意删去。根式化简还需保留符号，如 $\sqrt{x^2y}=|x|\sqrt y$ 在 $y\ge0$ 时成立；直接写成 $x\sqrt y$ 会在 $x<0,y>0$ 时出错。

有理化应服从目的。表达式 $\sqrt{x+1}-\sqrt x$ 在 $x\ge0$ 时等于 $1/(\sqrt{x+1}+\sqrt x)$；后者既显示它为正，也便于比较大小。这里分母严格为正，因此转换没有丢失取值。

\section{例题、习题与学习检查}
\begin{example}
求 $f(x)=x^3-2x^2-x+2$ 的全部实根，并求 $f(x)$ 除以 $x-3$ 的余式。
\end{example}
\begin{solve}
分组得 $f(x)=x^2(x-2)-(x-2)=(x-2)(x-1)(x+1)$，故根为 $2,1,-1$。余式由代入得到 $f(3)=27-18-3+2=8$。求根与求余式使用不同结构，不必为了余式完整做一次除法。
\end{solve}
\begin{enumerate}
\item 若 $x^2+ax+b=(x+2)^2$ 恒成立，求 $a,b$。答：$a=4,b=4$。
\item 化简 $(x^2-4)/(x^2+x-6)$ 并保留条件。答：$(x+2)/(x+3)$，$x\ne2,-3$。
\item 一个次数至多为 $3$ 的多项式在 $0,1,2,3$ 处均取值 $5$，证明它恒等于 $5$。提示：研究它与常数 $5$ 的差。
\end{enumerate}

